\chapter{Cardiac Rehabilitation}

\section*{Definition and Overview}
Cardiac rehabilitation (cardiac rehab) is a structured programme of exercise, education, and support for people who have had a heart attack, heart surgery, angioplasty, or who live with heart failure. Its aim is to help people recover physically and emotionally, reduce the risk of further heart problems, and return to normal activities with confidence. A typical programme combines supervised exercise sessions, information about heart-healthy living, and practical help with medication, smoking cessation, diet, and stress management. It is an evidence-based standard of care recommended internationally for almost everyone who has a cardiac event.

\section*{Epidemiology}
Heart disease is one of the leading causes of death and disability worldwide, so cardiac rehabilitation is offered to a very large group of people. It is recommended for almost everyone who has had a heart attack, a stent or bypass operation, valve surgery, or heart failure. Despite this, fewer than half of eligible people take part, with lower uptake among women, older people, and those from disadvantaged backgrounds. Participation also varies widely between countries and regions, and improving attendance is recognised as a major opportunity to reduce the burden of recurrent heart disease.

\section*{Aetiology and Pathophysiology}
The need for cardiac rehabilitation arises from the underlying heart condition and the effects of the event itself:
\begin{itemize}
    \item \textbf{Coronary artery disease:} the fatty narrowing of the heart's arteries that causes heart attacks and angina
    \item \textbf{Physical deconditioning:} bed rest and inactivity after a heart event cause muscle wasting, reduced fitness, and lowered exercise tolerance
    \item \textbf{Psychological effects:} anxiety, fear of exertion, depression, and reduced confidence are common after a cardiac event
    \item \textbf{Ongoing risk factors:} smoking, high blood pressure, high cholesterol, diabetes, overweight, and inactivity continue to drive disease if left unaddressed
    \item \textbf{Cardiac remodelling:} after a heart attack the heart may enlarge and weaken, and structured exercise can help counter this process
\end{itemize}

\section*{Clinical Features}
The programme is designed for people recovering from heart problems, whose typical symptoms and concerns include:
\begin{itemize}
    \item Reduced exercise tolerance and easy tiredness
    \item Breathlessness or chest discomfort on exertion
    \item Swelling of the ankles in heart failure
    \item Fear of overdoing it or of triggering another heart event
    \item Low mood, anxiety, or poor sleep after the event
    \item Uncertainty about when it is safe to return to work, driving, and sexual activity
\end{itemize}

\section*{Differential Diagnosis}
Rehabilitation is not a diagnosis but a treatment programme, so the main question is whether someone is ready and safe to participate. This is assessed against:
\begin{itemize}
    \item Worsening heart failure or unstable angina that needs treatment first
    \item Dangerous arrhythmias that need investigation before exercise is started
    \item Significant valve disease awaiting surgery
    \item Other health problems such as severe arthritis, frailty, or cognitive impairment that affect how the programme is delivered
    \item Symptoms such as breathlessness that may reflect lung disease rather than the heart
\end{itemize}

\section*{When to Seek Medical Care}
During cardiac rehabilitation, certain symptoms should always be reported before exercising or continuing a session:
\begin{itemize}
    \item Chest pain, pressure, or heaviness during or after exercise
    \item Severe breathlessness, or feeling faint, dizzy, or light-headed
    \item Palpitations or a very fast or irregular heartbeat
    \item Unusual swelling, rapid weight gain, or worsening symptoms between sessions
\end{itemize}
\textbf{Call 000 (or local emergency services) for:}
\begin{itemize}
    \item Chest pain lasting more than a few minutes or not relieved by rest
    \item Collapse or loss of consciousness
    \item Severe breathlessness at rest
\end{itemize}

\section*{Diagnosis and Investigations}
Before starting, each participant is assessed so that the programme can be tailored to their needs:
\begin{itemize}
    \item \textbf{Medical review:} the heart event, current medications, and any other health conditions
    \item \textbf{Exercise testing:} an ECG, often during a supervised stress test, to check the heart's response to exertion and set safe exercise limits
    \item \textbf{Blood tests:} cholesterol, blood sugar, and kidney function to guide risk-factor treatment
    \item \textbf{Blood pressure and weight:} measured at baseline and throughout the programme
    \item \textbf{Screening questionnaires:} for anxiety, depression, and readiness to change lifestyle habits
    \item \textbf{Ongoing monitoring:} heart rate, blood pressure, and symptoms are checked during supervised exercise sessions
\end{itemize}

\section*{Management}
A comprehensive cardiac rehabilitation programme has several components:
\begin{itemize}
    \item \textbf{Supervised exercise:} gradually increasing aerobic activity such as walking, cycling, or rowing, usually two to three times a week, plus strengthening exercises
    \item \textbf{Education:} understanding the condition, medicines, warning symptoms, and how to exercise safely
    \item \textbf{Dietary advice:} a heart-healthy eating pattern, weight management, and reduction of salt and saturated fat
    \item \textbf{Smoking support:} counselling and nicotine-replacement therapies for people who smoke
    \item \textbf{Medication review:} ensuring that heart-protective medicines are taken correctly
    \item \textbf{Psychological support:} help with anxiety and depression, and advice on stress management
    \item \textbf{Return to life:} practical guidance on work, driving, sport, and sex
    \item \textbf{Home-based options:} supported programmes delivered by phone, online, or with a paper workbook for those unable to attend in person
\end{itemize}

\section*{Complications}
Cardiac rehabilitation is generally very safe because exercise is supervised and individualised:
\begin{itemize}
    \item Chest discomfort or breathlessness if exercise is too intense
    \item Muscle aches, sprains, or joint pain from new activity
    \item Very rarely, a serious heart rhythm problem or heart event during exercise, the risk of which is minimised by pre-assessment and monitoring
    \item Emotional difficulties surfacing as people confront what has happened to them
    \item Falling behind or dropping out if the programme is not suited to a person's needs or transport is a barrier
\end{itemize}

\section*{Prognosis and Outcomes}
People who complete cardiac rehabilitation generally have a better outlook: fewer further heart events, fewer hospital readmissions, and improved survival compared with those who do not take part. Many regain fitness and confidence and return successfully to work and daily life. The benefits are seen across all age groups and in both men and women. Outcomes are best when rehabilitation is started soon after the event, completed in full, and followed by lasting changes in lifestyle.

\section*{Prevention and Self-Care}
\begin{itemize}
    \item Attend every session; the benefits grow with full participation
    \item Build up daily activity gradually between sessions, such as regular walking
    \item Follow the heart-healthy eating and weight-management advice you are given
    \item Stop smoking and keep alcohol within recommended limits
    \item Take prescribed medicines exactly as directed
    \item Learn your warning symptoms and know when to seek urgent help
    \item Keep active after the programme ends and seek a maintenance group if one is available
    \item Ask for support if you are struggling emotionally or with motivation
\end{itemize}

\section*{Sources and Further Reading}
The following sources provide authoritative, up-to-date information on cardiac rehabilitation and its benefits.
\begin{itemize}
    \item NHS. \textit{Cardiac rehabilitation overview.} https://www.nhs.uk/conditions/cardiac-rehabilitation/
    \item British Heart Foundation. \textit{Cardiac rehabilitation and recovery.} https://www.bhf.org.uk/
    \item National Heart, Lung, and Blood Institute. \textit{Cardiac rehabilitation.} https://www.nhlbi.nih.gov/health/cardiac-rehabilitation
    \item World Health Organization. \textit{Cardiovascular disease prevention and rehabilitation.} https://www.who.int/
    \item Heart Foundation Australia. \textit{Heart recovery and cardiac rehabilitation.} https://www.heartfoundation.org.au/
    \item Mayo Clinic. \textit{Cardiac rehabilitation: what to expect.} https://www.mayoclinic.org/tests-procedures/cardiac-rehabilitation/
\end{itemize}
